% Generated draft: requires mathematical and source review.
\documentclass[11pt]{article}
\usepackage[margin=25mm]{geometry}
\usepackage{fontspec}
\setmainfont{Latin Modern Roman}
\usepackage{amsmath,amssymb,mathtools,unicode-math}
\setmathfont{Latin Modern Math}
\usepackage{enumitem,needspace,xurl,hyperref}
\setlength{\parindent}{0pt}
\setlength{\parskip}{5pt}
\setlength{\emergencystretch}{3em}
\newcommand{\N}{\mathbb{N}}
\newcommand{\Z}{\mathbb{Z}}
\newcommand{\Q}{\mathbb{Q}}
\newcommand{\R}{\mathbb{R}}
\newcommand{\C}{\mathbb{C}}
\newcommand{\F}{\mathbb{F}}
\newcommand{\A}{\mathbb{A}}
\newcommand{\PP}{\mathbb{P}}
\newcommand{\E}{\mathbb{E}}
\newcommand{\eps}{\varepsilon}
\newcommand{\dd}{\,\mathrm{d}}
\begin{document}
\section*{UnsolvedMath ID: 701324}
\textbf{Source problem code:} catalogue entry
\subsection{Definitions and mathematical statement}
A variety \(\mathbf V\) of groups is the class of all groups satisfying a specified collection of group identities, with group homomorphisms as morphisms. For an infinite set \(X\), let \(\operatorname{Sym}(X)\) be the group of all bijections \(X\to X\), under composition. An embedding means an injective group homomorphism.

For \(A\in\mathbf V\), write \(A\amalg_{\mathbf V}A\) for the coproduct of two copies of \(A\) in \(\mathbf V\). Thus it is a group \(C\in\mathbf V\) equipped with homomorphisms \(\iota_1,\iota_2:A\to C\) such that, for every \(B\in\mathbf V\) and every pair of homomorphisms \(f_1,f_2:A\to B\), there is a unique homomorphism \(f:C\to B\) satisfying \(f\circ\iota_j=f_j\) for \(j=1,2\).

\textbf{Question (G. M. Bergman, Kourovka Notebook 17.22).} For every variety \(\mathbf V\) of groups, every infinite set \(X\), and every group \(A\in\mathbf V\), does the following implication hold?
\[
\bigl(\exists\text{ an injective homomorphism }A\to\operatorname{Sym}(X)\bigr)
\quad\Longrightarrow\quad
\bigl(\exists\text{ an injective homomorphism }A\amalg_{\mathbf V}A\to\operatorname{Sym}(X)\bigr).
\]
The same set \(X\) is used on both sides; no compatibility with a chosen embedding of \(A\) is requested.

The supplied statement attributes an affirmative answer for the variety of all groups to N. G. de Bruijn. The supplied registry labels the general problem open while recording an explicitly unverified AI claim of a negative answer. That claim is not treated here as an established resolution.
\subsection{Short English statement}
If a group \(A\) in a variety \(\mathbf V\) embeds into the full symmetric group on an infinite set \(X\), must the coproduct in \(\mathbf V\) of two copies of \(A\) embed into that same symmetric group? The source attributes the affirmative unrestricted-group case to N. G. de Bruijn; its reported general negative answer is explicitly unverified.
\subsection{Sources}
\begin{itemize}
\item Supplied local UnsolvedMath registry record, permanent ID \(701324\), exact source code \texttt{catalogue entry}, title \emph{Kourovka Notebook Problem 17.22}. Proposed by G. M. Bergman in 2010.
\item \emph{The Kourovka Notebook (Unsolved Problems in Group Theory)}, problem 17.22, 17th issue (2010); the supplied record also identifies the 21st issue edition (2026), \texttt{arXiv:1401.0300v48}. The collection description names E. I. Khukhro and V. D. Mazurov as editors.
\item N. G. de Bruijn: the supplied problem statement attributes to him the affirmative result when \(\mathbf V\) is the variety of all groups. No bibliographic reference for that result is supplied.
\item Achyuth Jayadevan, \emph{Preprint}, 2026, \url{https://zenodo.org/records/22909835}. The supplied discussion explicitly labels its reported negative answer an ``unverified AI claim'' and its proposed solution an ``unverified AI solution''; no independent verification is asserted here.
\item Supplied status review dated \texttt{2026-10-16}: canonical status \texttt{open}, evidence classification \texttt{unrefereed\_claim}. Its notes report a \(\#\) marker for the unverified AI solution claim in the October 2026 21st issue, \texttt{arXiv:1401.0300v48}. This is supplied dated curation, not an independently verified literature assessment.
\end{itemize}
\end{document}
